\section{HDR Floating-Point Representation}
\label{sec:hdr-float}

Deep-zoom Mandelbrot rendering places unusual demands on numerical
representation. Orbit values may span an extreme dynamic range over the course
of iteration, while small relative differences between nearby orbits must be
tracked accurately. Standard floating-point formats are poorly matched to this
combination of requirements: fixed-precision mantissas limit relative accuracy
at large magnitudes, while arbitrary-precision formats incur prohibitive
computational cost when used per pixel.

To address this gap, the implementation employs a custom \emph{high dynamic range (HDR)}
floating-point representation that explicitly separates scale and precision
(\cref{fig:type-hierarchy}).
This representation is designed to preserve relative accuracy across a wide
range of magnitudes while remaining efficient on GPUs.

\subsection{Decoupling scale and precision}

In standard IEEE floating-point arithmetic, each value is represented as a
mantissa multiplied by a power of two, with both components stored implicitly in
a single word. While convenient, this tightly couples scale and precision: as
values grow in magnitude, fewer bits remain available to represent small
relative differences.

The HDR representation instead stores values in the form
\begin{equation}
x = m \cdot 2^{e},
\end{equation}
where:
\begin{itemize}
\item \(m\) is a fixed-precision mantissa stored in standard floating-point or
      double-double format, and
\item \(e\) is an explicitly managed integer exponent.
\end{itemize}

By storing the exponent separately, the mantissa can remain centered near unit
magnitude regardless of the overall scale of the value. This separation preserves relative
precision even when the absolute value becomes extremely large or small.

\subsection{Normalization and arithmetic}

HDR arithmetic maintains the invariant that the mantissa remains within a fixed
normalized range. After each arithmetic operation, the mantissa is renormalized
and any excess scaling is folded into the exponent. Conceptually, this mirrors
the behavior of floating-point normalization, but with explicit control over the
process.

Arithmetic operations are performed as follows:
\begin{itemize}
\item \textbf{Multiplication} combines mantissas and adds exponents.
\item \textbf{Addition and subtraction} align exponents explicitly before
      operating on mantissas.
\item \textbf{Renormalization} ensures that mantissas remain well-scaled and
      numerically stable.
\end{itemize}

Because mantissas are stored using fixed-precision types, these operations map
naturally to GPU hardware while avoiding the overhead of general-purpose
arbitrary-precision arithmetic.

\subsection{HDR complex numbers}

Mandelbrot iteration operates on complex numbers. In the HDR formulation, each
complex value
\[
z = x + i y
\]
is represented by storing separate HDR values for the real and imaginary
components. Both components share the same conceptual structure but may have
independent exponents.

This representation allows the magnitude of \(z\) to grow or shrink over many
orders of magnitude while maintaining accurate relative phase and magnitude
information, which is critical for both escape testing and perturbation-based
methods.

\subsection{Role of HDR in perturbation and reference orbits}

The HDR representation plays a complementary role to perturbation
(\cref{sec:perturbation-concept}).  It extends range, not significand precision.
The application-authoritative CPU recurrence uses MPIR, and the CUDA
reference-orbit recurrence uses fixed-width \code{HpSharkFloat}.  Both emit
reduced HDR samples for downstream work.

HDR arithmetic is used primarily for:
\begin{itemize}
\item storing reduced reference-orbit samples and parameters,
\item perturbation and linear-approximation arithmetic whose values require wide
      range but only the selected mantissa precision, and
\item GPU periodicity, escape, radius, and derivative calculations around the
      \code{HpSharkFloat} orbit state.
\end{itemize}

HDR therefore enables the mixed-precision structure used by perturbation
rendering; it does not replace the MPIR or \code{HpSharkFloat} recurrence.

\subsection{Advantages over conventional representations}

Compared to standard floating-point formats, HDR arithmetic offers:
\begin{itemize}
\item substantially increased dynamic range,
\item stable relative precision independent of magnitude,
\item predictable and controllable numerical behavior.
\end{itemize}

Compared to arbitrary-precision libraries, it avoids dynamic memory allocation,
irregular control flow, and large per-operation overheads, making it well suited
for parallel execution on GPUs.

\subsection{Summary}

The HDR floating-point representation provides a numerically robust and
computationally efficient foundation for deep-zoom Mandelbrot rendering. By
explicitly separating scale and precision, it bridges the gap between standard
floating-point arithmetic and full arbitrary-precision methods. In combination
with perturbation, it enables accurate, high-performance rendering across
extreme zoom levels while remaining compatible with GPU execution models.

\begin{figure}[!htbp]
\centering
\begin{tikzpicture}[
    >=Stealth,
    typebox/.style={draw, rounded corners=2pt, minimum height=0.55cm,
                    minimum width=1.8cm, font=\scriptsize\ttfamily,
                    fill=blue!8, align=center},
    bitlabel/.style={font=\tiny, gray},
]

% ═══ Part (a): precision hierarchy ═══
\node[font=\scriptsize\bfseries] at (1.6, 5.8) {(a) Precision hierarchy};

% Row 0: native types
\node[typebox] (float) at (0, 0) {float};
\node[bitlabel, right=1pt] at (float.east) {24\,b};
\node[typebox] (double) at (3.2, 0) {double};
\node[bitlabel, right=1pt] at (double.east) {53\,b};

% Row 1: 2-component
\node[typebox] (dblflt) at (0, 1.4) {dblflt};
\node[bitlabel, right=1pt] at (dblflt.east) {${\sim}$48\,b};
\node[typebox] (dbldbl) at (3.2, 1.4) {dbldbl};
\node[bitlabel, right=1pt] at (dbldbl.east) {${\sim}$104\,b};

% Row 2: 4-component
\node[typebox] (gqf) at (0, 2.8) {GQF};
\node[bitlabel, right=1pt] at (gqf.east) {${\sim}$96\,b};
\node[typebox] (gqd) at (3.2, 2.8) {GQD};
\node[bitlabel, right=1pt] at (gqd.east) {${\sim}$212\,b};

% Row 3: reference-orbit arithmetic
\node[typebox, fill=gray!12] (mpir) at (0, 4.2) {MPIR};
\node[bitlabel, right=1pt] at (mpir.east) {runtime-selected};
\node[typebox, fill=gray!12] (hpshark) at (3.2, 4.2) {HpSharkFloat};
\node[bitlabel, right=1pt] at (hpshark.east) {fixed specialization};

% Arrows: increasing precision
\draw[->, blue!40, thick] (float) -- (dblflt);
\draw[->, blue!40, thick] (double) -- (dbldbl);
\draw[->, blue!40, thick] (dblflt) -- (gqf);
\draw[->, blue!40, thick] (dbldbl) -- (gqd);
\draw[->, blue!40, thick] (gqf.north) -- (mpir.south);
\draw[->, blue!40, thick] (gqd.north) -- (hpshark.south);

% Column labels
\node[font=\tiny, gray, rotate=90, anchor=south] at (-1.3, 1.4)
  {fp32-based};
\node[font=\tiny, gray, rotate=90, anchor=south] at (5.4, 1.4)
  {fp64-based};

% HDRFloat<T> wrapper
\begin{scope}[on background layer]
  \node[draw, dashed, rounded corners=6pt, orange!70!black, thick,
        inner xsep=16pt, inner ysep=8pt,
        fit=(float)(double)(gqf)(gqd)] (hdrbox) {};
  \node[font=\scriptsize, orange!70!black, anchor=south east]
    at (hdrbox.north east) {\code{HDRFloat<T>} wrapper};
\end{scope}

% Precision arrow
\draw[->, thick, gray!50] (-1.8, -0.3) -- (-1.8, 4.5);
\node[font=\tiny, gray, rotate=90, anchor=south] at (-2.1, 2.1)
  {increasing precision};

% ═══ Part (b): HDR representation ═══
\begin{scope}[shift={(7.5, 0.5)}]
  \node[font=\scriptsize\bfseries] at (2.0, 5.3)
    {(b) HDR representation};

  % IEEE 754 double
  \node[font=\scriptsize, anchor=east] at (-0.1, 3.65)
    {IEEE 754 \code{double}:};

  \draw[fill=gray!20] (0, 3.3) rectangle (0.4, 4.0);
  \node[font=\tiny, rotate=90] at (0.2, 3.65) {S};
  \draw[fill=blue!15] (0.4, 3.3) rectangle (1.5, 4.0);
  \node[font=\tiny] at (0.95, 3.65) {exp};
  \draw[fill=blue!8] (1.5, 3.3) rectangle (4.0, 4.0);
  \node[font=\tiny] at (2.75, 3.65) {mantissa};

  \draw[decorate, decoration={brace, amplitude=4pt, mirror}, thick]
    (0, 3.15) -- (4.0, 3.15)
    node[midway, below=5pt, font=\tiny]
      {64~bits, tightly coupled};

  % HDRFloat<T>
  \node[font=\scriptsize, anchor=east] at (-0.1, 1.45)
    {\code{HDRFloat<T>}:};

  \draw[fill=orange!20, draw=orange!70!black, thick]
    (0, 1.1) rectangle (1.5, 1.8);
  \node[font=\tiny] at (0.75, 1.45) {int32 exp};

  \node[font=\scriptsize] at (1.75, 1.45) {$+$};

  \draw[fill=blue!12, draw=blue!60!black, thick]
    (2.0, 1.1) rectangle (4.0, 1.8);
  \node[font=\tiny] at (3.0, 1.45) {\code{T} mantissa};

  \draw[decorate, decoration={brace, amplitude=4pt, mirror},
        thick, orange!70!black]
    (0, 0.95) -- (1.5, 0.95)
    node[midway, below=5pt, font=\tiny, orange!70!black] {range};
  \draw[decorate, decoration={brace, amplitude=4pt, mirror},
        thick, blue!60!black]
    (2.0, 0.95) -- (4.0, 0.95)
    node[midway, below=5pt, font=\tiny, blue!60!black] {precision};

  \node[font=\tiny, text width=3.5cm, align=center, gray] at (2.0, 0.1)
    {decoupled: scale and precision\\vary independently};
\end{scope}

\end{tikzpicture}
\caption[Numeric type hierarchy and HDR representation]{%
  Numeric type hierarchy and HDR representation.
  (a)~The precision hierarchy of floating-point types available in
  \FractalShark{}, from native IEEE~754 types (bottom) through
  multi-component expansion types to the two reference-orbit arithmetic roles
  at the top: runtime-selected MPIR and fixed-width \code{HpSharkFloat}
  specializations.  The dashed outline indicates the mantissa types used by the
  implemented \code{HDRFloat<T>} instantiations; it is not a claim that every
  device type is wrapped.
  (b)~Comparison of the standard IEEE~754 \code{double} layout, where
  exponent and mantissa share a fixed bit budget, with the
  \code{HDRFloat<T>} representation, which stores an explicit integer
  exponent separately from a fixed-precision mantissa of type~\code{T}.}
\label{fig:type-hierarchy}
\end{figure}
